\section{Reflection and the plane normalization}
\label{sec:plane}

Throughout this section the law is the free plane pair obtained in
Proposition~\ref{prop:local-mixing}, with the fixed weight
$x\in[1/\sqrt2,1]$.  Only its original nearest-neighbor interaction is used
for reflection; no reflection symmetry of the auxiliary cell variables is
needed.  Write $\xi_u=H(u)-H(0)$ for its height increments.  If
$f\in C_c^\infty(\R^2)$ has integral zero, set
\[
 H_\delta(f)=\sum_{u\in\T}w_{\delta,f}(u)H(u),\qquad
 w_{\delta,f}(u)=\int_{\delta(u+V)}f(z)\dd z,
\]
where $V$ is the closed unit-mesh hexagonal face centered at zero, with
boundaries ignored in integrals.  The faces $u+V$ tile the plane.  In
particular, $\sum_u w_{\delta,f}(u)=\int f=0$, so the statistic is
independent of the additive lift.  All these finite sums are centered by
height reflection.  We will show that their covariances converge to a
positive multiple of the plane Green form, with the same constant along
every mesh sequence.  We then use this covariance to show that a Laplace
test supported strictly on one side of a reflection line has covariance
with its reflected test tending to zero.  Reflection positivity will
interpret this as a vanishing reflection seminorm, which
Section~\ref{sec:pins} will use to transfer the Laplace identity through
separated pinning conditions.

Use Euclidean coordinates, and put
\[
 e_1=(1,0),\qquad e_2=(1/2,\sqrt3/2),\qquad
 P=\frac{2\pi}{\sqrt3},\qquad
 \T^*=\{(2\pi m,P(2n-m)):m,n\in\Z\}.
\]
The reciprocal torus is $\widehat\T=\R^2/\T^*$, with characters
$u\mapsto e^{ip\cdot u}$.  Small neighborhoods of its origin always use
their Euclidean representatives.  We use
$\widehat f(p)=\int e^{-ip\cdot z}f(z)\dd z$ for the continuum Fourier
transform.

The covariance statement below is the output of the spectral argument.
Its coefficient will be fixed by a full scalar cutoff limit, after the
possible subsequential spectra have been identified.  The proof is completed
in the final subsection; the reflection-null insertion then follows from it.

\begin{proposition}[Plane covariance and normalization]
\label{prop:plane-covariance}
There is a constant $v=v(x)\in(0,\infty)$, depending only on the fixed
weight $x$, such that for every pair of real zero-integral tests
$f,g\in C_c^\infty(\R^2)$,
\begin{align}
 \lim_{\delta\downarrow0}\E[H_\delta(f)H_\delta(g)]
 &=\frac{v}{(2\pi)^2}
       \int_{\R^2}\frac{\widehat f(p)\overline{\widehat g(p)}}{|p|^2}
                  \dd p                                      \label{plane:eq:covariance}\\
 &=\frac{v}{2\pi}\iint_{\R^2\times\R^2}
          f(z)\log\frac1{|z-w|}\,g(w)\dd z\dd w.              \notag
\end{align}
The conclusion also holds for zero-integral tests converging in
$C_c^\infty$ with a common compact support.
\end{proposition}

\subsection{Reflection positivity and normal translation}

The row-conditioning proof below is the site-reflection construction
of \cite[Section~3, example~5]{FILS1978} specialized to the present edge law.
We include the normal-translation argument because positivity of that
translation, not only reflection positivity of the field, is needed
for the nonnegative Cauchy mixture.

\begin{proposition}[Positive normal transfer]
\label{prop:reflection-transfer}
Let $\ell$ be any row of lattice sites parallel to one of the three
unoriented nearest-neighbor directions, and choose one of its closed
half-planes.  The plane pair is reflection positive across $\ell$.
The assertion remains valid after adjoining any finite collection of fields
\begin{equation}
 \label{plane:eq:phases}
 Z_s(u)=U_s\exp\{is(H(u)-H(o))\},
\end{equation}
where the $U_s$ are independent unit-circle Haar variables, independent of
the pair, and $o$ is any anchor.

More precisely, let $\theta_\ell$ act on observables by spatial reflection,
and take local functions supported in the chosen half-plane, whose
height differences are evaluated along paths in that half-plane.  Then
\[
 \langle F,G\rangle_\ell
     =\E\big[\overline{\theta_\ell F}\,G\big]
\]
is a positive semidefinite Hermitian form.  In its quotient completion,
translation farther into the half-plane by the pure normal vector of
length $\sqrt3$ is a positive self-adjoint contraction $S$.  Unit
translation along the row is a commuting unitary $U$.  At mesh $\delta$
the normal step is $\sqrt3\delta$.
\end{proposition}

\begin{proof}
Take first a reflection-symmetric finite free volume, chosen so that its
intersection with the row is connected and contains the observables under
consideration.  No edge of the triangular lattice jumps across a row of
sites.  Conditional on the row spins, the two strict sides therefore have
independent, reflected-identical laws: every original edge weight belongs
to a single side together with its row boundary.  Row spins determine
relative row heights.  If phases are present, condition also on their
values at one row anchor.  On each side the other phases are determined
from those anchor phases and the increments on that side.  Their
conditional Haar law at the row anchor is independent of the spins, even
if the original anchor $o$ was elsewhere.  Consequently
\[
 \E\big[\overline{\theta_\ell F}\,F\big]
   =\E\left[\left|\E[F\mid\text{row spins and row phases}]\right|^2\right]
   \ge0.
\]
The plane limit gives the same inequality for local functions.  Ordinary
$L^2$ approximation is also allowed, since
$|\langle F,G\rangle_\ell|\le\|F\|_2\|G\|_2$.
Changing the anchor or translating the configuration multiplies each
$U_s$ by a unit complex number determined by the increments.  Conditional
Haar measure is unchanged.  This proves stationarity and all the lattice
reflection symmetries also for the augmented law.

For the operator statement take $\ell=\{y=0\}$ and the upper half-plane.
A site has coordinates $(m+n/2,n\sqrt3/2)$.  Reflections in $y=0$ and
$y=\sqrt3/2$ send its integer coordinates respectively to
$(m+n,-n)$ and $(m+n-1,2-n)$.  Thus both are site-row reflections, and their
composition gives the pure normal translation $a=(0,\sqrt3)$.
Let $T$ denote translation of observables by $a$.  Stationarity and
$\theta_\ell T=T^{-1}\theta_\ell$ imply
\[
 \langle TF,G\rangle_\ell=\langle F,TG\rangle_\ell.
\]
Furthermore $TF$ is supported above the halfway row $y=\sqrt3/2$, and
reflection of $TF$ in that row is $\theta_\ell F$.  Reflection positivity
at that row gives $\langle F,TF\rangle_\ell\ge0$.

Boundedness must be proved before taking the quotient.  Put
$a_n=\langle T^nF,T^nF\rangle_\ell^{1/2}$.  Symmetry and reflection
Cauchy--Schwarz yield
\[
 a_n^2=\langle T^{n-1}F,T^{n+1}F\rangle_\ell
       \le a_{n-1}a_{n+1}\quad(n\ge1),
 \qquad a_n\le\|F\|_2.
\]
If $a_0=0$, the first inequality gives $a_1=0$.  Otherwise an increase
$a_1>a_0$ would, by log-convexity, force geometric growth, contradicting
the second inequality.  Hence $\|TF\|_\ell\le\|F\|_\ell$.
The operator descends to the quotient and extends to a bounded symmetric
positive operator on its completion, hence to a positive self-adjoint
contraction $S$.  Tangential translation preserves the form and has an
inverse with the same property.  It therefore gives a unitary $U$;
$U,U^{-1}$ commute with $S$.  Rotation by $\pi/3$ and $2\pi/3$, and
translation of the base row, give exactly the same construction in all
three directions.
\end{proof}

\subsection{The increment spectrum and its Cauchy representation}

\begin{lemma}[Punctured spectral measure]
\label{plane:lem:spectrum}
There is a positive, locally finite measure $C$ on
$\widehat\T\setminus\{0\}$, invariant under the lattice rotations and
reflections, such that every finitely supported real array $(b_u)$ with
$\sum b_u=0$ satisfies
\begin{equation}
 \label{plane:eq:spectrum}
 \E\left|\sum_u b_uH(u)\right|^2
   =\int\left|\sum_u b_ue^{ip\cdot u}\right|^2 C(\dd p).
\end{equation}
There is no invariant random linear tilt term.  Moreover, for some
$r_0>0$ and $K<\infty$,
\begin{equation}
 \label{plane:eq:annuli}
 C\{r\le|p|\le2r\}\le K\quad(0<r<r_0),
 \qquad
 \int_{0<|p|<r_0}|p|^t C(\dd p)<\infty\quad(t>0).
\end{equation}
\end{lemma}

\begin{proof}
Let $V_u$ be ordinary unitary translation on the plane probability space.
Lemma~\ref{spectra:unitary}, applied to the two basis translations,
constructs their scalar covariance measures $\mu_{a,b}(\dd p)$ on
$\widehat\T$.  Write $P_0$ for the common invariant projection.
All increments belong to $L^2$, since
$|\xi_u|$ is bounded by the graph distance from zero to $u$.  The cocycle
identity gives
\[
 \xi_{u+v}=\xi_v+V_v\xi_u,
 \qquad (V_v-I)\xi_u=(V_u-I)\xi_v.
\]
Put $d_u(p)=e^{ip\cdot u}-1$.  The polynomial covariance identity in
Lemma~\ref{spectra:unitary} transfers the cocycle relation to the scalar
measures.  On a patch where $d_v\ne0$, it gives
$\mu_{\xi_u,\xi_w}=\overline{d_u}d_w
\mu_{\xi_v,\xi_v}/|d_v|^2$.  Division is made first on compact
subpatches where the denominator is bounded below.  Thus define there
\[
 C(\dd p)=
 \frac{\mu_{\xi_v,\xi_v}(\dd p)}
      {|e^{ip\cdot v}-1|^2}.
\]
The definitions agree on overlaps.  The two choices $v=e_1,e_2$ cover
the punctured torus; on every compact subset their denominators are
bounded below on a finite cover.  This proves local finiteness and
\eqref{plane:eq:spectrum} for the part orthogonal to $P_0$.

The remaining projected cocycle is additive:
$P_0\xi_u=A\cdot u$ for some random vector $A\in L^2$.
For a smooth compactly supported $f$ of integral zero, use the exact
face weights for $R^{-2}f(z/R)$ at unit mesh.  Their first lattice moment is
\[
 \sum_u u\int_{u+V}R^{-2}f(z/R)\dd z
       =R\int zf(z)\dd z+O(1),
\]
where the deterministic error is bounded because $|u-z|$ is bounded
on $u+V$.  Orthogonal projection and Proposition~\ref{prop:plane-moments}
bound the $L^2$ norm of the dot product of $A$ with this vector uniformly
in $R$.  Divide by $R$ and let $R\to\infty$.  Choosing two such $f$ with
independent first moments proves $A=0$.  Orthogonality rules out any
cancellation between this tilt and other spectral components.
Uniqueness of the local construction of $C$ gives its stated symmetries.

For the annular estimate choose a real smooth bump $b$ with integral one
and support so small that $\widehat b$ has no zero on $|p|\le2$.
The tests $f_1=\partial_1b$ and $f_2=\partial_2b$ have integral zero, and
$|\widehat f_1|^2+|\widehat f_2|^2$ is bounded below on
$1\le|p|\le2$.  Let $b^{(r,j)}_u$ be the unit-face weights for
$r^2f_j(rz)$.  Their absolute sums are bounded and their largest
absolute values are $O(r^2)$; their supports have radius $O(r^{-1})$.
For $1\le|p|\le2$,
\[
 \sum_u b^{(r,j)}_ue^{irp\cdot u}\longrightarrow\widehat f_j(-p)
 \quad\text{uniformly as }r\downarrow0.
\]
Indeed replacing the face center by a point in its face changes the
exponential by $O(r)$ uniformly here.  Apply
Proposition~\ref{prop:plane-moments} and \eqref{plane:eq:spectrum} to
the two weights.  Their squared transforms bound below a positive
constant on $r\le|p|\le2r$, proving the first assertion in
\eqref{plane:eq:annuli} for small $r$.  Local finiteness handles any
remaining compact range of radii.  Summation over dyadic annuli gives
the second assertion.
\end{proof}

Restrict translations to the rectangular sublattice generated by $e_1$
and $(0,\sqrt3)$.  Its reciprocal torus has centered rectangle
\[
 \mathcal R=[-\pi,\pi)\times[-P/2,P/2).
\]
The projection $\pi:\widehat\T\to\mathcal R$ is a two-sheeted covering.
The preimages of zero are zero and the nonzero alias $a_*=(0,P)$ on
$\widehat\T$.  Let $C'=\pi_*C$, with the value of $C$ at the omitted
origin set to zero when interpreting this pushforward.
For $0<s<\infty$ define the probability measure $Q_s$ on the normal
circle by
\begin{equation}
 \label{plane:eq:cauchy}
 Q_s(\dd l)=q_s(l)\dd l,\qquad
 q_s(l)=\sum_{n\in\Z}\frac{s}{\pi\{s^2+(l+nP)^2\}}.
\end{equation}
Here and below an integral over $l$ uses the centered period.
Set $Q_0=\delta_0$ and let $Q_\infty$ be uniform probability on that circle.

\begin{lemma}[Positive Cauchy mixture and exceptional sets]
\label{plane:lem:mixture}
On $k\ne0$ in the rectangular torus there is a positive measure
$\nu(\dd k,\dd s)$ such that
\begin{equation}
 \label{plane:eq:mixture}
 C'(\dd k,\dd l)=\int_{[0,\infty]}Q_s(\dd l)\,\nu(\dd k,\dd s).
\end{equation}
The measure $\nu$ is finite on horizontal bands bounded away from zero,
and, for every sufficiently small fixed $\kappa>0$,
\begin{equation}
 \label{plane:eq:weighted-nu}
 \int_{0<|k|<\kappa}|k|\,\nu(\dd k,\dd s)<\infty.
\end{equation}
The physical punctured line $k=0$ near zero has zero $C$ mass, and
$C\{a_*\}=0$.  These conclusions hold before any rescaling limit is taken.
The analogous assertions hold in each of the three rotated charts.
\end{lemma}

\begin{proof}
Take $F=\xi_{e_1}$, an observable on the horizontal reflection row.
Apply the scalar construction of Lemma~\ref{spectra:transfer} to the
commuting $U,S$ from Proposition~\ref{prop:reflection-transfer}.
It gives a finite positive measure $\gamma(\dd k,\dd\lambda)$ on the
unit circle times $[0,1]$ such that, for every $m\in\Z$ and $j\ge0$,
\[
 \E[F\,V_{me_1+j(0,\sqrt3)}F]
   =\langle F,U^mS^jF\rangle_\ell
   =\int e^{imk}\lambda^j\,\gamma(\dd k,\dd\lambda).
\]
Reflection in the base row gives the same expression with $|j|$ for
negative $j$.  Put $\lambda=e^{-\sqrt3s}$, interpreting $\lambda=0$ as
$s=\infty$.  The $j$th Fourier coefficient of $Q_s$ is
$e^{-\sqrt3s|j|}$, including both endpoints.  Uniqueness of Fourier
coefficients of finite measures on the rectangular torus identifies
the ordinary spectral measure of $F$ with this mixture.
By \eqref{plane:eq:spectrum} that measure is
$|e^{ik}-1|^2 C'(\dd k,\dd l)$.  Dividing by this strictly positive
factor on $k\ne0$ gives \eqref{plane:eq:mixture} and local finiteness
of $\nu$ there.
Since $Q_s$ has total mass one, the left side of
\eqref{plane:eq:weighted-nu} equals
$\int_{0<|k|<\kappa}|k|\,C'(\dd k,\dd l)$.
Near the physical origin this is finite by \eqref{plane:eq:annuli} and
$|k|\le|p|$; its remaining preimage has finite mass by local finiteness
on the physical torus away from zero.  This includes a neighborhood of
$a_*$ and does not require any estimate on horizontal bands near zero.

For the physical line take a compact segment of $k=0$ in a small
punctured disk and rotate it through $\pi/3$.  Its coordinates become
$(k',l')=(-\sqrt3l/2,l/2)$: both are nonzero, and the image is a
nonvertical graph away from $l'=0$.  For each fixed $k'$ its normal
section is a singleton outside zero.  Every $Q_s$ assigns zero mass
to this singleton: $Q_s$ is nonatomic for $s>0$, including $s=\infty$,
while $Q_0$ is supported at zero.  Applying \eqref{plane:eq:mixture}
and then positivity of the pushforward gives zero mass to the rotated
segment, hence to the original segment.  Countably many segments cover
the punctured line.
For the alias, rotation sends $(0,P)$ to $(-\pi,P/2)$.
On the rectangular torus this point has nonzero tangential and nonzero
normal coordinates.  The same mixture argument gives zero mass to that
point.  Projection can only add mass, so invariance under rotation yields
$C\{a_*\}=0$.  In particular this argument does not identify the two
boundary representatives of the rectangle on the physical torus.
Rotating the entire construction proves the assertions for the other
row directions.
\end{proof}

\subsection{A summable dispersion defect}

The unperiodized Cauchy kernel with $s=|k|$ has density
$|k|/\{\pi(k^2+l^2)\}$.  We next show that only mixing mass near this
relation survives at small frequencies.  The three mirror directions
will then determine the tangential measure as well.  The key estimate
comes from averaging a fourth angular harmonic, whose rotational average
vanishes but whose Cauchy average measures the deviation from $s=|k|$.

For centered coordinates write
\[
 \chi(k,l)=\cos(4\arg(k+il))
   =1-\frac{8k^2l^2}{(k^2+l^2)^2},\qquad (k,l)\ne0,
\]
and put
\[
 D(k,s)=\left(\frac{|k|-s}{|k|+s}\right)^2\quad(s<\infty),
 \qquad D(k,\infty)=1.
\]

\begin{lemma}[Angular identity and summable defect]
\label{plane:lem:defect}
Uniformly for $s\in[0,\infty]$ and sufficiently small $k\ne0$,
\begin{equation}
 \label{plane:eq:poisson}
 \int\chi(k,l)Q_s(\dd l)=D(k,s)+O(|k|).
\end{equation}
For some fixed $\kappa>0$,
\begin{equation}
 \label{plane:eq:defect}
 \int_{0<|k|<\kappa}D(k,s)\,\nu(\dd k,\dd s)<\infty.
\end{equation}
There is also $K_1<\infty$ such that, for all sufficiently small $r>0$,
\begin{equation}
 \label{plane:eq:bands}
 \nu\{r\le|k|\le2r\}\le K_1.
\end{equation}
\end{lemma}

\begin{proof}
Put $a=|k|$.  The analytic function
$((a+iz)/(a-iz))^2$ is bounded in the upper half-plane, has boundary
real part $\chi(a,l)$, and has real value $((a-s)/(a+s))^2$ at $z=is$.
The Poisson formula therefore gives the exact line identity
\begin{equation}
 \label{plane:eq:line-poisson}
 \int_\R\chi(a,l)\frac{s\dd l}{\pi(s^2+l^2)}
     =\left(\frac{a-s}{a+s}\right)^2\quad(0<s<\infty).
\end{equation}
Here the Poisson formula applies to a bounded harmonic function with
continuous boundary values, so no boundary growth condition is missing.
To estimate periodization, let
$f_a(l)=1-\chi(a,l)=8a^2l^2/(a^2+l^2)^2$.  Direct integration gives
\[
 \int_\R f_a(l)\dd l=4\pi a,
 \qquad 0\le f_a(l)\le\frac{8a^2}{l^2}\quad(l\ne0).
\]
If $p_s(l)=s/\{\pi(s^2+l^2)\}$, then
\[
 M_P:=\sup_{s>0,\ |l|\le P/2}\sum_{n\ne0}p_s(l+nP)<\infty.
\]
Indeed for $s\le P$ use $|l+nP|\ge P(|n|-1/2)$ and sum the
square tail; for $s>P$ the $O(s/P)$ terms with $|n|\le2s/P+1$
are each at most $1/(\pi s)$ and the remaining terms have a summable
$s/(nP)^2$ bound.  The difference between the periodized and line
integrals of $f_a$ is at most
\[
 4\pi M_Pa+
 \int_{|l|>P/2}f_a(l)p_s(l)\dd l
 \le4\pi M_Pa+32a^2/P^2.
\]
This proves \eqref{plane:eq:poisson} uniformly for finite positive $s$.
For $s=0$ it is exact; for $s=\infty$ use the uniform density and
$\int f_a\le4\pi a$.

Integrate \eqref{plane:eq:poisson} over
$\epsilon<|k|<\kappa$.  Its error is bounded independently of
$\epsilon$ by \eqref{plane:eq:weighted-nu}.
Replacing the projected integral by the physical integral inside
$|p|<\kappa$ changes it by a bounded amount: the remaining physical
region, including the other sheet, has finite mass away from zero.
Thus it suffices to bound above
\[
 \int_{\epsilon<|p|<\kappa,\ |k|>\epsilon}
                \chi(k,l)\,C(\dd p).
\]
The integral over the full radial annulus is exactly zero, since
$\sum_{j=0}^5\cos(4(\arg p+j\pi/3))=0$ and $C$ is rotation invariant.
In the deleted cap $|k|\le\epsilon$, the negative part of $\chi$ is
confined to $|l|<(1+\sqrt2)|k|$, hence to
$\epsilon<|p|<K_2\epsilon$ for an absolute constant $K_2$.
A fixed number of annular bounds in \eqref{plane:eq:annuli} controls
that negative part.  Deleting the positive part can only lower the
integral.  We have obtained an upper bound independent of $\epsilon$.
Since $D\ge0$, monotone convergence proves \eqref{plane:eq:defect}.
This use of a negative cap requires only annular bounds, not the
quantitative endcap estimate proved below.

On $r\le|k|\le2r$, if $s\notin[r/2,4r]$ then $D(k,s)\ge1/9$;
\eqref{plane:eq:defect} bounds the mass of this part.
For $s\in[r/2,4r]$, the central, unperiodized part of $Q_s$ gives
\[
 Q_s\{|l|\le2r\}\ge\frac2\pi\arctan(1/2)>0
\]
when $r$ is small.  Equation~\eqref{plane:eq:mixture} therefore bounds
the remaining $\nu$ mass by a constant times
$C'\{r\le|k|\le2r,\ |l|\le2r\}$.  The physical-origin preimage
lies in a fixed number of annuli of scale $r$, and the alias preimage
has bounded finite mass.  This proves \eqref{plane:eq:bands}.
\end{proof}

\begin{lemma}[Endcaps]
\label{lem:endcaps}
After decreasing $r_0$ if necessary, there is $K_3<\infty$ such that
\begin{equation}
 \label{plane:eq:endcaps}
 C\{b<|p|\le2b,\ |k|\le h\}\le K_3\frac hb
 \qquad(0<b<r_0,\quad0<h<b/10).
\end{equation}
The same bound holds in each rotated chart.  In particular, rescaling
cannot create mass on one of the tangential zero lines inside a compact
punctured annulus.
\end{lemma}

\begin{proof}
Rotate the indicated cap through $\pi/3$.  Both new coordinates $k',l'$
have magnitude comparable to $b$.  At fixed $k'$, the inequality
$|k|\le h$ restricts $l'$ to an interval of length $4h/\sqrt3$.
For example $0.8b\le|k'|\le1.8b$ and
$0.4b\le|l'|\le1.1b$ suffice when $h<b/10$.
The central Cauchy density is bounded there, uniformly in $s>0$, by
$1/(2\pi|l'|)\le K/b$; its folded tails are bounded by $M_P$.
The uniform endpoint obeys the same bound, and the $s=0$ atom misses
the interval.  The relevant horizontal band has bounded $\nu$ mass
by \eqref{plane:eq:bands}, enlarging by finitely many dyadic bands
if necessary.  Apply \eqref{plane:eq:mixture} and note that projection
adds nonnegative mass to the rotated physical cap.  This proves
\eqref{plane:eq:endcaps}.  Applying the bound with $b=\delta b_0$
and $h=\delta h_0$, then sending $h_0\downarrow0$, proves the last
assertion for every vague limit; continuous cutoff majorants justify
passage to such limits.
\end{proof}

\subsection{The full spectral scaling limit}

For $\delta>0$ rescale a fixed small physical neighborhood of zero by
$p\mapsto p/\delta$, and denote its image measure by $C_\delta$.
On every fixed compact subset of $\R^2\setminus\{0\}$ this is
well defined for small $\delta$.  Equation~\eqref{plane:eq:annuli}
gives local boundedness and hence subsequential vague compactness.

\begin{lemma}[Form of every subsequential spectrum]
\label{plane:lem:subsequences}
Every vague subsequential limit is
\begin{equation}
 \label{plane:eq:subsequence}
 M(\dd p)=c\frac{\dd p}{|p|^2}
\end{equation}
for some finite $c\ge0$, at this point possibly depending on the
subsequence.  This identity holds on the entire punctured plane,
including the tangential zero lines.
\end{lemma}

\begin{proof}
Fix $0<a<b<\infty$ and work on the rescaled horizontal band
$a\le|k|/\delta\le b$.  Rescale the mixing variables to
$(k/\delta,s/\delta)$.  Their masses are bounded by finitely many
instances of \eqref{plane:eq:bands}.  Their integral of $D$ tends to
zero by \eqref{plane:eq:defect}, since the unscaled bands shrink to
zero while excluding it.  On this band the scale-invariant function
$D$ is bounded below away from zero whenever $s/\delta$ stays outside
any fixed neighborhood of $|k|/\delta$; this includes $s/\delta=\infty$.
Thus every limiting mixing measure is concentrated on the graph
$s=|k|$ in the rescaled coordinates.  More explicitly, mass outside
$|s/|k|-1|<\varepsilon$ is bounded by a constant depending on
$\varepsilon$ times the integral of $D$, and tends to zero.
This also gives tightness in the rescaled $s$ coordinate.

Uniformly for $a/2\le t\le2b$ and $l$ in a compact set,
\[
 \delta q_{\delta t}(\delta l)
     =\frac{t}{\pi(t^2+l^2)}+O(\delta),
\]
by the folded-tail bound in the proof of
Lemma~\ref{plane:lem:defect}.  Consequently every further extraction
of the bounded mixing measures gives, on $k\ne0$,
\begin{equation}
 \label{plane:eq:limiting-mixture}
 M(\dd k,\dd l)=
     \frac{|k|}{\pi(k^2+l^2)}\,\eta(\dd k)\dd l
\end{equation}
for some locally finite positive measure $\eta$ in the tangential
coordinate.  A diagonal extraction over compact horizontal subbands
makes these descriptions consistent.  The other physical preimage of
a bounded rescaled set shrinks to $a_*$; its mass tends to zero by
local finiteness there and $C\{a_*\}=0$, already proved in
Lemma~\ref{plane:lem:mixture}.  Thus \eqref{plane:eq:limiting-mixture}
indeed describes the physical measure, with no alias contribution.

Multiplying \eqref{plane:eq:limiting-mixture} by $|p|^2$ shows that
the normal distributional derivative of $|p|^2M$ vanishes on $k\ne0$.
No density assumption on $\eta$ is being made.  Repeating in the
other two row directions, at every nonzero point at least two
of the three tangential coordinates are nonzero.  The corresponding
normal directions are independent.  On a neighborhood of that point
the two independent directional derivatives of $|p|^2M$ vanish.
A distribution with this property is constant there: convolve on a
smaller ball with a smooth approximate identity, observe that both
derivatives of each convolution vanish, and pass to the limit.
The constants agree on overlaps, and the punctured plane is connected.
It follows that $|p|^2M=c\dd p$ there.  This argument covers the
axes by the other charts; Lemma~\ref{lem:endcaps} also directly
excludes concentration on them.  Positivity and the annular bound
give $0\le c<\infty$.
\end{proof}

The annular bound and the preceding lemma alone would not imply that
$c$ is unique.  We next construct a scalar limit before selecting a
spectral subsequence.

\begin{lemma}[A cutoff identity determining the coefficient]
\label{plane:lem:unique}
There is a single $c\ge0$ such that $C_\delta$ converges vaguely to
$c\dd p/|p|^2$ as $\delta\downarrow0$.
\end{lemma}

\begin{proof}
Choose an even smooth $q_0$ with $0\le q_0\le1$, equal to one on
$[-1/2,1/2]$ and zero outside $(-1,1)$.  Let $Q$ be a smooth radial
cutoff supported in a sufficiently small physical disk and equal to
one near zero.  For $h>0$ put
\begin{equation}
 \label{plane:eq:cutoff}
 J(h)=\int Q(p)\chi(k,l)
       [q_0(|p|/h)-q_0(k/h)]\,C(\dd p).
\end{equation}
Each integrand vanishes near zero and has compact support in the
physical chart, so this is a finite integral.
We first show that $J(h)$ has a full limit as $h\downarrow0$.
Split its bracket as
\[
 [q_0(|p|/h)-1]+[1-q_0(k/h)].
\]
The first term has integral zero by sixfold rotation, since its other
factors except $\chi$ are radial and it vanishes near zero.
For the second term replace the physical cutoff $Q$ by the rectangular
strip $0<|k|<\kappa$ in \eqref{plane:eq:mixture}.  The resulting error
has a limit by dominated convergence: in a sufficiently small physical
neighborhood of zero the cutoffs and centered coordinates agree;
all remaining contributions come from a finite measure away from
that origin, including the alias neighborhood.  Their multipliers
are bounded and converge pointwise as $h\downarrow0$.  On $k=0$
the factor $1-q_0(k/h)$ is identically zero.

For the strip contribution set
$A(k,s)=\int\chi(k,l)Q_s(\dd l)$.  Equations
\eqref{plane:eq:poisson}, \eqref{plane:eq:defect}, and
\eqref{plane:eq:weighted-nu} give
\[
 |A(k,s)|\le D(k,s)+K|k|,
 \qquad
 \int_{0<|k|<\kappa}|A(k,s)|\,\nu(\dd k,\dd s)<\infty.
\]
Dominated convergence applied to
$\int(1-q_0(k/h))A(k,s)\nu(\dd k,\dd s)$ proves that the strip
contribution has a limit.  Thus
\begin{equation}
 \label{plane:eq:cutoff-full-limit}
 L:=\lim_{h\downarrow0}J(h)
\end{equation}
exists, using only the summable defect and finite cutoff errors,
without any assumption about a spectral scaling limit.

Now take any sequence along which $C_h\to c\dd p/|p|^2$.
In the rescaled variable the bracket in \eqref{plane:eq:cutoff} is
$q_0(|p|)-q_0(k)$, which vanishes on $|p|\le1/2$.
On a dyadic annulus $R<|p|\le2R$ with $R>10$, it is supported on
$|k|\le1$.  Equation~\eqref{plane:eq:endcaps}, used at physical
scales $b=hR$ and horizontal width $h$, bounds that mass by $K/R$.
Summing dyadic annuli gives a uniform $O(1/R)$ tail, up to the edge
of the fixed support of $Q$.  Vague convergence on the remaining
compact annulus therefore gives
\begin{align}
 L&=c\int_{\R^2}\frac{\chi(k,l)}{|p|^2}
              [q_0(|p|)-q_0(k)]\dd p \notag\\
  &=c\int_0^{2\pi}\cos(4u)
       \int_0^\infty[q_0(r)-q_0(r|\cos u|)]\frac{\dd r}{r}\dd u
    =-\frac\pi2c.                       \label{plane:eq:cutoff-value}
\end{align}
For completeness, for $t>0$ the inner radial integral with $t=|\cos u|$
equals $\log t$: integrate first from $\varepsilon$ to $M$, change
variables in the second term, then use that $q_0=1$ near zero and
vanishes for large argument.  Its absolute integral is bounded by
$K(1+|\log t|)$, which is integrable in $u$, justifying Fubini.
The angular value can be checked without a Fourier-series limit:
\[
 \int_0^{2\pi}\cos(4u)\log|\cos u|\dd u
 =4\int_0^{\pi/2}\cos(4u)\log(\cos u)\dd u
 =4\int_0^{\pi/2}\sin^2u\cos(2u)\dd u
 =-\frac\pi2.
\]
The middle equality follows by integration by parts; the boundary
term $\sin(4u)\log(\cos u)$ vanishes at both ends.
Since $L$ was defined by a full limit, \eqref{plane:eq:cutoff-value}
fixes the same $c$ for every subsequence.  Subsequence compactness
then gives full vague convergence.
\end{proof}

\subsection{Euclidean area covariance and reflection-null insertions}

\begin{proof}[Proof of Proposition~\ref{prop:plane-covariance}]
For the exact area weights define
$A_{\delta,f}(p)=\sum_uw_{\delta,f}(u)e^{ip\cdot u}$.
On every fixed compact set of $\zeta$,
\begin{equation}
 \label{plane:eq:transform-limit}
 A_{\delta,f}(\delta\zeta)\longrightarrow\widehat f(-\zeta)
 \quad\text{uniformly}.
\end{equation}
Indeed the defining sum is the integral of $f(z)$ times the
exponential at its face center.  Replacing that center by $z$
changes the exponential by $O(\delta|\zeta|)$ on each face.
The total integral is exact, so no lattice-cell area factor remains.

We record bounds that also control frequencies outside a fixed
rescaled annulus.  For every integer $N\ge1$, in a fixed small
physical neighborhood of zero,
\begin{equation}
 \label{plane:eq:transform-bounds}
 |A_{\delta,f}(p)|\le K_N
       \min\left\{\frac{|p|}{\delta},
                    (1+|p|/\delta)^{-N}\right\},
\end{equation}
and on the complement of that neighborhood it is $O_N(\delta^N)$.
The first bound follows from exact zero total and
$\sum_u|w_{\delta,f}(u)|\,|u|=O(\delta^{-1})$.
For the other bounds, let $\Delta_e$ denote a unit lattice difference
in a basis direction.  Translation of faces gives
\[
 \sum_u|\Delta_e^Nw_{\delta,f}(u)|
   \le\int_{\R^2}|\Delta_{\delta e}^N f(z)|\dd z
   \le\delta^N\|\partial_e^N f\|_{L^1}.
\]
Discrete summation by parts therefore bounds
$|1-e^{ip\cdot e}|^N|A_{\delta,f}(p)|$ by this last quantity.
Near zero at least one of the two basis multipliers is at least
$a|p|$; away from zero on the reciprocal torus at least one is bounded
below on a finite compact cover.  Together with
$|A_{\delta,f}|\le\|f\|_1$, these estimates prove
\eqref{plane:eq:transform-bounds}.  Their constants are uniform for
a bounded smooth family with a common compact support.

Equations~\eqref{plane:eq:spectrum} and
\eqref{plane:eq:annuli} now allow passage to the full Fourier integral.
For example the contribution of $|p|<\epsilon\delta$ to a variance
is $O(\epsilon^2)$ by the first bound in
\eqref{plane:eq:transform-bounds} and dyadic summation.
The contribution of $R\delta<|p|<r_0$ is $O_N(R^{-2N})$;
the remaining compact part of the torus contributes $O_N(\delta^{2N})$
by local finiteness.  On $\epsilon\delta\le|p|\le R\delta$ use
Lemma~\ref{plane:lem:unique} and \eqref{plane:eq:transform-limit}.
Let first $\delta\downarrow0$, then $\epsilon\downarrow0$ and
$R\to\infty$.  This gives the variance limit
$c\int|\widehat f(p)|^2|p|^{-2}\dd p$; polarization gives the
bilinear formula.  The same estimates prove the assertion for
smoothly converging tests.

The distribution $(2\pi)^{-1}\log(1/|z|)$ is the fundamental
solution of $-\Delta$ on the plane.  Its ambiguity by a constant is
irrelevant to zero-integral tests.  Fourier inversion with the stated
convention thus gives the second line of \eqref{plane:eq:covariance}
and the normalization
\begin{equation}
 \label{plane:eq:normalization}
 v=(2\pi)^2c,\qquad \sigma(x)=\sqrt v.
\end{equation}
Finiteness of $c$ was already proved.  Lemma~\ref{lem:nondegeneracy}
supplies two fixed smooth nonnegative unit averages on separated
disks whose difference $f$ has
$\inf_{\delta\text{ small}}\E H_\delta(f)^2>0$.
The covariance limit therefore rules out $c=0$.
All constructions have used only the fixed plane law at weight $x$;
there is no domain or approximation parameter in $v$.
\end{proof}

\begin{proposition}[Reflection-null Laplace insertion]
\label{prop:reflection-null}
Fix one of the three row directions.  Let $\ell_\delta$ be site rows
of $\delta\T$ converging to a line $\ell$ in that direction, and
choose the same side of each row.  Suppose that real
$\varphi_\delta\to\varphi$ in $C_c^\infty(\R^2)$, with common
compact support lying on that side at a positive distance from
$\ell_\delta$ for all sufficiently small $\delta$.  Put
$X_\delta=H_\delta(\Delta\varphi_\delta)$, using exact area weights.
Then
\begin{equation}
 \label{plane:eq:reflection-null}
 \|X_\delta\|_{\ell_\delta}^2
   =\E[(\theta_{\ell_\delta}X_\delta)X_\delta]
       \longrightarrow0,
 \qquad \sup_{\delta\text{ small}}\|X_\delta\|_{L^2}<\infty.
\end{equation}
This holds also in every finitely phase-augmented reflection space of
Proposition~\ref{prop:reflection-transfer}.  In particular, if
$Y_\delta$ is supported on the same side and has bounded reflection
norm, then
$\E[(\theta_{\ell_\delta}Y_\delta)X_\delta]\to0$.
\end{proposition}

\begin{proof}
All Laplace tests have integral zero.  Reflection in a site row
permutes the hexagonal faces exactly, so
\[
 \theta_{\ell_\delta}X_\delta
   =H_\delta\bigl(\Delta(\varphi_\delta\circ\theta_{\ell_\delta})\bigr)
\]
as zero-total increment observables.  Their supports, including the
finitely many faces meeting the test supports, lie strictly on
opposite sides once $\delta$ is small.  The tests and their reflections
converge smoothly on fixed compact sets.  By
Proposition~\ref{prop:plane-covariance}, the limit of their covariance is
\[
 \frac{v}{(2\pi)^2}\int |p|^2\widehat\varphi(p)
             \overline{\widehat{\varphi\circ\theta_\ell}(p)}\dd p
    =v\int\nabla\varphi(z)\cdot
                   \nabla(\varphi\circ\theta_\ell)(z)\dd z=0,
\]
because the two supports are disjoint.  Reflection positivity makes
the finite-mesh quantity a squared seminorm.  The ordinary second
moments are bounded by Proposition~\ref{prop:plane-moments}, or by
Proposition~\ref{prop:plane-covariance}.  Adjoining independent Haar
variables does not change the reflection norm of the phase-independent
$X_\delta$.  Finally apply reflection Cauchy--Schwarz.  Only the
reflection seminorm vanishes; an ordinary variance need not do so.
\end{proof}
